\section*{第 \ref{ch:b2:fluid-kinematics} 章——流体运动学}

\begin{solution}{exo:b2:fluid-kinematics:1}
(a) $\ell/L \sim 7 \times 10^{-8}$：连续介质。(b) $0.07$：临界——
连续介质定律开始失效（壁面滑移）。(c) $\ell = 70 \times 10^{-9}
/3 \times 10^{-7} = \qty{0.23}{m}$，$\ell/L = 0.23$：不是连续介质；气体
是由独立分子组成的“分子雨”（自由分子区）。
\end{solution}

\begin{solution}{exo:b2:fluid-kinematics:2}
$a = v\,\dd v/\dd x = v_0^2(1 + x/L)/L$（从 $v_0^2/L$ 增至 $2v_0^2/L$）。
$\dd x/\dd t = v_0(1 + x/L)$：$t = (L/v_0)\ln2 = 0.69\,L/v_0$，介于
$L/2v_0$ 与 $L/v_0$ 之间。
\end{solution}

\begin{solution}{exo:b2:fluid-kinematics:3}
$D_V = \int_0^Rv_{\max}(1 - r^2/R^2)\,2\pi r\,\dd r = \tfrac12\pi R^2v_{\max}$；
$\langle v\rangle = v_{\max}/2 = \qty{0.20}{m/s}$；$D_V = \qty{1.6e-5}{m^3/s} =
\qty{16}{mL/s}$。
\end{solution}

\begin{solution}{exo:b2:fluid-kinematics:4}
(a) $\operatorname{div} = 0$，$\vect\omega = \vect 0$：不可压缩、
无旋。(b) $\operatorname{div} = 0$，$\omega_z = 2\Omega$：刚体
转动。(c) $\operatorname{div} = 2a \ne 0$，无旋（平面
源）。(d) $\operatorname{div} = 0$，$\omega_z = -k$：剪切流。(e)
$\operatorname{div} = b$：仅当 $b = 0$ 时不可压缩；对
每个 $b$ 都无旋。
\end{solution}

\begin{solution}{exo:b2:fluid-kinematics:5}
$D_V = \qty{0.20}{L/s}$。软管 $S = \qty{1.13}{cm^2}$：$v = \qty{1.8}{m/s}$；喷管
$S = \qty{7.1}{mm^2}$：$v = \qty{28}{m/s}$。平均加速度 $\Delta(v^2/2)/L
= (28^2 - 1.8^2)/0.08 \approx \qty{1e4}{m/s^2} \approx 1000\,g$。时间 $\approx
\int\dd x/v$：约几毫秒（线性速度分布时为 \qty{4}{ms}）。
\end{solution}

\begin{solution}{exo:b2:fluid-kinematics:6}
(a) $\dd y/\dd x = (V/U)\cos\omega t$：斜率为
$(V/U)\cos\omega t$ 的直线，彼此平行，且随时间倾斜。(b) $x = Ut$，$y =
(V/\omega)\sin\omega t$：正弦曲线 $y = (V/\omega)\sin(\omega x/U)$。(c) 在
$t = 0$ 时，\omterm{def:b2:fluid-kinematics:euler}{流线}与横轴成 \qty{45}{\degree}，\omterm{def:b2:fluid-kinematics:euler}{迹线}
从原点出发时与横轴成 \qty{45}{\degree}；在 $t = \pi/2\omega$ 时\omterm{def:b2:fluid-kinematics:euler}{流线}
水平，而\omterm{def:b2:fluid-kinematics:euler}{迹线}为固定的正弦曲线。两者不同，是因为
流动非定常：\omterm{def:b2:fluid-kinematics:euler}{流线}是某一时刻的快照，\omterm{def:b2:fluid-kinematics:euler}{迹线}则是质点的历史。
(d) $\vect a = \partial_t\vect v = (0, -V\omega\sin\omega t)$——该场是
均匀的，因此对流项为零。
\end{solution}

\begin{solution}{exo:b2:fluid-kinematics:7}
(a) $D_V = 50 \times 2 \times 1.2 = \qty{120}{m^3/s}$。(b) $120/60 = \qty{2.0}{m/s}$。
(c) $150/150 = \qty{1.0}{m/s}$。(d) 若 $\rho \propto 1/T$，且压强 $P$ 不变，则
密度减半；$\rho vS$ 守恒，速率加倍。
\end{solution}

\begin{solution}{exo:b2:fluid-kinematics:8}
(a) $\omega_z = \frac1r\frac{\dd(rv_\theta)}{\dd r}$：内部为 $2\Omega$，外部
为 $0$。(b) 内部 $\Gamma = 2\pi r\cdot\Omega r = 2\Omega\cdot\pi r^2$，即
\omterm{def:b2:fluid-kinematics:vorticity}{涡量}的通量；外部 $\Gamma = 2\pi\Omega a^2$，为常数，即
整个涡核的通量。(c) $\Omega = v_{\max}/a = \qty{1.4}{rad/s}$；$\Gamma = 2\pi av_{\max}
= \qty{2.2e4}{m^2/s}$；在 \qty{500}{m} 处：$\Omega a^2/r = \qty{7}{m/s}$。(d) 涡核
像刚体一样转动；外部流动无旋：质点
绕轴运动，自身却不发生自转。\omterm{def:b2:fluid-kinematics:vorticity}{涡量}——
即旋转——只存在于涡核中。
\end{solution}

\begin{solution}{exo:b2:fluid-kinematics:9}
(a) $\varphi = \tfrac12k(x^2 - y^2)$，$\Delta\varphi = k - k = 0$。(b) $\dd x/
kx = \dd y/(-ky)$：$xy = $ 常数；$\partial_y\psi = kx = v_x$，$-\partial_x\psi =
-ky = v_y$。(c) $\vect a = (v_x\partial_x + v_y\partial_y)\vect v = k^2(x, y) =
\operatorname{\vect{grad}}\bigl(\tfrac12k^2(x^2 + y^2)\bigr) = \operatorname{\vect{grad}}
(v^2/2)$：在定常\omterm{def:b2:fluid-kinematics:vorticity}{无旋流动}中，加速度是
单位质量动能的梯度；它指向远离
驻点的方向——质点沿 $y$ 方向接近壁面时减速，
沿 $x$ 方向远离时加速。(d) $\dot x = kx$，$\dot y = -ky$：
$x = x_0\eu^{kt}$，$y = y_0\eu^{-kt}$（$xy$ 为常数）。
\end{solution}

\begin{solution}{exo:b2:fluid-kinematics:10}
(a) $\dd x/\dd t = x/(t + \tau)$ 给出 $x = x_0(1 + t/\tau)$，$\ddot x = 0$；
\omterm{thm:b2:fluid-kinematics:material}{随体导数}：$\partial_tv + v\partial_xv = -x/(t + \tau)^2 + x/(t + \tau)^2
= 0$。(b) $\partial_t\rho + \partial_x(\rho v) = \dot\rho + \rho/(t + \tau) = 0$：$\rho =
\rho_0\tau/(t + \tau)$。(c) 两质点间距按 $(1 + t/\tau)$ 增长，
而它们之间的质量固定：$\rho \propto 1/(1 + t/\tau)$。(d) $v =
Hx$，$H = 1/(t + \tau)$：哈勃定律，自由膨胀（无
加速度），密度与尺度因子成反比下降。
\end{solution}

\begin{solution}{exo:b2:fluid-kinematics:11}
(a) $\vect v = \operatorname{\vect{grad}}\varphi = (U + qx/2\pi r^2, qy/2\pi r^2)$；
源的 $v_r = q/2\pi r$，$\operatorname{div} = \frac1r\dd(rv_r)/\dd r
= 0$，且与 $\theta$ 无关，因此 $\omega_z = 0$；这类场的叠加
仍是不可压缩、无旋的。(b) 在负 $x$ 轴上，
$U - q/2\pi r = 0$：$x_s = -q/2\pi U$。(c) 在驻点 $\theta =
\pi$，$y = 0$ 处：$\psi = q/2$；\omterm{def:b2:fluid-kinematics:euler}{流线} $Ur\sin\theta + q\theta/2\pi = q/2$，
即 $r\sin\theta = (q/2\pi U)(\pi - \theta)$——一条从 $x_s$ 出发、
向下游张开的曲线。(d) $\theta \to 0$：$y \to q/2U$；$\theta \to 2\pi$（下
支）：$y \to -q/2U$。该\omterm{def:b2:fluid-kinematics:euler}{流线}内侧的流体来自
源；其外侧则是均匀来流绕过一个
渐近宽度为 $q/U$ 的半体。
\end{solution}

\begin{solution}{exo:b2:fluid-kinematics:12}
(a) $\operatorname{div}\vect v = \frac1r\partial_r(-q/2\pi) = 0$；$\omega_z =
\frac1r\partial_r(\Gamma/2\pi) = 0$。(b) 由 $\dd r/v_r = r\,\dd\theta/v_\theta$ 可得
$\dd r/r = -(q/\Gamma)\dd\theta$。(c) $a_r = v_r\partial_rv_r - v_\theta^2/r =
-(q^2 + \Gamma^2)/4\pi^2r^3 = -v^2/r$；$a_\theta = v_r\partial_rv_\theta + v_rv_\theta/r
= 0$。(d) $\dd r/\dd t = -q/2\pi r$：$t = \pi(r_0^2 - a^2)/q$；圈数
$= (\Gamma/2\pi q)\ln(r_0/a)$。(e) $t = \qty{520}{s} \approx \qty{9}{min}$；$10$
圈。
\end{solution}

\begin{solution}{pb:b2:fluid-kinematics:1}
\textbf{1.} 穿过半径为 $r$ 的圆柱面：$2\pi rh|v_r| = D_V$，于是
$v_r = -q/2\pi r$, $q = D_V/h = \qty{1.5e-3}{m^2/s}$.

\textbf{2.} $rv_\theta = 0.5 \times 0.01 = \qty{5e-3}{m^2/s}$: $v_\theta = \Gamma/2\pi r$
其中 $\Gamma = \qty{3.1e-2}{m^2/s}$。

\textbf{3.} $\operatorname{div}\vect v = \frac1r\partial_r(rv_r) = 0$, $\omega_z =
\frac1r\partial_r(rv_\theta) = 0$.

\textbf{4.} $\dd r/r = -(q/\Gamma)\dd\theta$：$r = r_0\eu^{-q\theta/\Gamma}$；圈数
$= (\Gamma/2\pi q)\ln(r_0/a) = 3.3 \times 3.2 \approx 11$.

\textbf{5.} $t = \pi(r_0^2 - a^2)/q = \qty{520}{s} \approx \qty{9}{min}$; $v_\theta(a) =
\qty{0.25}{m/s}$.

\textbf{6.} $\vect a = -(v^2/r)\vect e_r$；当 $r = a$ 时，$v^2 = 0.25^2 + 0.012^2$，
$a = \qty{3.1}{m/s^2} \approx 0.3g$.

\textbf{7.} $\Omega_\oplus\sin\lambda\,r = 5 \times 10^{-5} \times 0.5 = \qty{2.5e-5}{m/s}$,
比搅拌作用小四百倍：半球效应
在普通浴缸中毫无影响（若在静置一天的水箱中则有影响）。

\textbf{8.} $v_\theta = \Omega r$ ($r < a$), $\Omega a^2/r$ ($r > a$); $\Omega = 80/60
= \qty{1.33}{rad/s}$, $\Gamma = 2\pi av_{\max} = \qty{3.0e4}{m^2/s}$.

\textbf{9.} $\omega_z = 2\Omega = \qty{2.7}{s^{-1}}$ 在核心区内，外部为 $0$；
$v_\theta$ 线性增至 \qty{80}{m/s}，随后按 $1/r$ 减小；$\omega_z$ 呈
阶跃状。

\textbf{10.} $a = v_{\max}^2/a = 6400/60 = \qty{107}{m/s^2} \approx 11g$，指向中心。

\textbf{11.} 外部\omterm{def:b2:fluid-kinematics:vorticity}{涡量}为零：一小块流体微团沿其圆周作\omterm{rem:b2:rigid-body-mechanics:motions}{平动}
而不绕自身转动——桨轮
保持其取向。在核心区（刚体旋转）中，它随
流体转动，每绕行一圈自转一次。

\textbf{12.} $E_{\text{核心区}} = \int_0^a\tfrac12\rho\Omega^2r^2\,2\pi r\,\dd r =
\tfrac14\pi\rho\Omega^2a^4 = \qty{2.2e7}{J/m}$.

\textbf{13.} $E_{\text{外部}} = \int_a^R\tfrac12\rho(\Omega a^2/r)^2\,2\pi r\,\dd r =
\pi\rho\Omega^2a^4\ln(R/a) = 8.7 \times 10^7 \times 4.4 = \qty{3.8e8}{J/m}$；被积函数
$\rho v^2r \propto 1/r$，故出现对数。

\textbf{14.} $\qty{4.0e8}{J/m} \times \qty{1000}{m} = \qty{4e11}{J}$：相当于约
\qty{100}{t} TNT 当量。

\textbf{15.} 流入量 $2\pi a \times 300 \times 5 = \qty{5.7e5}{m^3/s}$ 穿过
面积 $\pi a^2 = \qty{1.1e4}{m^2}$：$v_z = \qty{50}{m/s}$。

\textbf{16.} $\omega = 2v_{\max}/a = \qty{2.5e-3}{s^{-1}} = 50f$; $\Gamma = 2\pi av_{\max}
= \qty{1.3e7}{m^2/s}$.

\textbf{17.} $r_0\cdot fr_0/2 = r(v_\theta + fr/2)$: $v_\theta = f(r_0^2 - r^2)/2r$;
\qty{60}{m/s} 对应 \qty{100}{km}，\qty{155}{m/s} 对应 \qty{40}{km}。

\textbf{18.} 在北半球，$f > 0$ 给出 $v_\theta > 0$：
从上方看为逆时针方向（气旋式）。在 \qty{500}{km} 尺度上，
行星自转是主要的初始旋转；在浴缸中
则是搅拌。

\textbf{19.} 与海面的摩擦带走角动量，而
压强梯度无法平衡
越转越快的环状气流所需的向心加速度：空气在到达轴线前
即在眼墙处上升，留下平静的台风眼。

\textbf{20.} $\rho\,2\pi rHv = 1.2 \times 2\pi \times 5 \times 10^5 \times 10^3 \times 5 =
\qty{1.9e10}{kg/s}$.

\textbf{21.} $\vect v(M, t) = -U\vect e_x + \vect v_0(M + Ut\,\vect e_x)$：图样
被平流输运，场在地面参考系中非定常。
\omterm{rem:b2:rigid-body-mechanics:motions}{平动}速度与涡旋叠加于 $v_\theta$ 指向
西的一侧——即西行轨迹的北（右手）侧。

\textbf{22.} $0.05 \times 10^7/3600 = \qty{139}{m^3/s}$; $v = 139/40 = \qty{3.5}{m/s}$.

\textbf{23.} $D_V = \qty{83}{cm^3/s}$：主动脉 \qty{0.33}{m/s}；毛细血管
\qty{0.33}{mm/s}；每毫米 \qty{3}{s}。

\textbf{24.} $S\,\dd h/\dd t = -s\sqrt{2gh}$: $\sqrt h = \sqrt{h_0} -
\frac sS\sqrt{g/2}\,t$，在 $T = (S/s)\sqrt{2h_0/g} = 10^4 \times 0.45 =
\qty{4500}{s} \approx \qty{75}{min}$ 时排空。

\textbf{25.} 浴缸：$r \sim \qty{0.5}{m}$，$v \sim \qty{1}{cm/s}$，$\Gamma =
\qty{3e-2}{m^2/s}$；龙卷风：\qty{60}{m}，\qty{80}{m/s}，$\omega = \qty{2.7}{s^{-1}}$，
$\Gamma = \qty{3e4}{m^2/s}$；气旋：\qty{40}{km}，\qty{50}{m/s}，$\omega =
\qty{2.5e-3}{s^{-1}}$，$\Gamma = \qty{1.3e7}{m^2/s}$。在各情形中，会聚流体的
单位质量角动量 $rv_\theta$ 守恒：
流入使其加速旋转。
\end{solution}
